\documentclass[11pt]{article}
\usepackage{amsmath,amssymb,amsthm}
\newtheorem{proposition}{Proposition}
\begin{document}
\title{Creation and annihilation parents: a nesting derivation of the existing second-order bounds}
\author{}\date{}\maketitle
This note derives the established bounds from structured conditional
parents. It does not assert a stronger unrestricted upper bound.

Let $\mathcal C\subset\operatorname{Sym}^N(\mathbb C^2)$ be a
$k$-dimensional second-order AC subspace with isometry $U$, and put
$j=N/2$, $c=j(j+1)/3$. Thus
\[
 U^\dagger J_aU=0,\qquad U^\dagger J_aJ_bU=c\delta_{ab}I_k.
\]
Use Schwinger bosons $a,b$ with total occupation $N=2j$,
$J_z=(a^\dagger a-b^\dagger b)/2$ and $J_+=a^\dagger b$.

\section*{Two structured parents}
Define creation and annihilation isometries, with columns grouped by
spinor label,
\[
 E_+=\frac{1}{\sqrt{j+1}}(a^\dagger U,\ b^\dagger U)
 :\mathbb C^2\otimes\mathbb C^k\longrightarrow V_{j+1/2},
\]
\[
 E_-=\frac{1}{\sqrt j}(aU,\ bU)
 :\mathbb C^2\otimes\mathbb C^k\longrightarrow V_{j-1/2}.
\]
The first-order AC equations imply $E_\pm^\dagger E_\pm=I_{2k}$.
The second-order equations give
\[
 E_+^\dagger(J_x',J_y',J_z')E_+
 =\alpha_+(S_x,S_y,S_z)\otimes I_k,
 \quad\alpha_+=\frac{N+3}{3},
\tag{1}
\]
\[
 E_-^\dagger(J_x',J_y',J_z')E_-
 =\alpha_-(S_x,-S_y,S_z)\otimes I_k,
 \quad\alpha_-=\frac{N-1}{3},
\tag{2}
\]
where $S_a=\sigma_a/2$. Conjugating the spinor factor of (2) by a
half-turn about $y$ turns the triple into $-\alpha_-(S_x,S_y,S_z)$.
The sign/orientation is essential: the annihilation parent carries the
dual spinor compression. Neither parent is a first-order AC subspace.
They instead have fully prescribed nonzero spin compressions.

For clarity, the diagonal $z$ entries follow directly from
\[
 aJ_z'a^\dagger=(j+J_z+1)(J_z+1/2),\qquad
 a^\dagger J_z'a=(j+J_z)(J_z-1/2),
\]
whose compressions, after normalization, are respectively
\[
 \frac{c+(j+1)/2}{j+1}=\frac{N+3}{6},\qquad
 \frac{c-j/2}{j}=\frac{N-1}{6}.
\]
The opposite spinor labels give their negatives. The nonzero raising
cross entries are supplied by
\[
 aJ_+'b^\dagger=(j+1)^2-J_z^2,\qquad
 b^\dagger J_+'a=j^2-J_z^2.
\]
The reverse cross entries are $J_+^2$ and compress to zero; diagonal
raising entries are linear combinations of $J_+$ and $J_+J_z$ and
also compress to zero. This proves (1) and (2) for arbitrary complex
AC2 subspaces, without a symmetry assumption.

\section*{Annihilation parent gives $4k\le N-2$}
Couple the physical spin $J=j-1/2=(N-1)/2$ to an auxiliary spin one,
and put $X=\mathbf J'\cdot\mathbf T$. On the $6k$-dimensional
structured parent times the auxiliary space, its compression is unitarily
similar to $-\alpha_-\mathbf S\cdot\mathbf T\otimes I_k$.
Its two eigenvalues are
\[
 -\alpha_-/2\quad\text{with multiplicity }4k,
 \qquad \alpha_-\quad\text{with multiplicity }2k.
\]
The ambient $X$ has eigenvalues and multiplicities
\[
\begin{array}{c|ccc}
L&J-1&J&J+1\\\hline
X_L&-(J+1)&-1&J\\
2L+1&N-2&N&N+2.
\end{array}
\]
For $N\ge8$, $-\alpha_-/2=-(N-1)/6<-1$.
Cauchy interlacing therefore forces the $4k$ compression eigenvalues
below $-1$ to fit into the ambient lowest eigenspace of dimension $N-2$:
\[
 \boxed{4k\le N-2\qquad(N\ge8).}
\]

\section*{Creation parent gives $5k\le N+5$}
Now couple physical spin $J=j+1/2=(N+1)/2$ to auxiliary spin $s=3/2$.
The compression of $X=\mathbf J'\cdot\mathbf T$ is
$\alpha_+\mathbf S\cdot\mathbf T\otimes I_k$. Its eigenvalues are
\[
 3\alpha_+/4=\frac{N+3}{4}\quad\text{with multiplicity }5k,
 \qquad -5\alpha_+/4\quad\text{with multiplicity }3k.
\]
For $N\ge2$, the two highest ambient eigenvalues are
\[
 X_{J+s}=\frac{3J}{2},\qquad X_{J+s-1}=\frac J2-\frac32=\frac{N-5}{4}.
\]
The positive compression eigenvalue is strictly above the second of these.
Only the ambient top total-spin sector can lie at or above it. Its
dimension is
\[
 2(J+s)+1=N+5.
\]
Interlacing therefore gives
\[
 \boxed{5k\le N+5.}
\]
Combined with the annihilation result, this recovers
\[
 K(N,2)\le\min\left\{\left\lfloor\frac{N-2}{4}\right\rfloor,
                         \left\lfloor\frac{N+5}{5}\right\rfloor\right\}
 \qquad(N\ge8).
\]
This is an alternative derivation of the existing inequalities, not an
improvement to their numerical values.

\section*{Other auxiliary spins and the limits of this calculation}
For an arbitrary auxiliary spin $s$, the creation-parent compression
has eigenvalues
\[
 \alpha_+s/2\ (\text{multiplicity }(2s+2)k),\qquad
 -\alpha_+(s+1)/2\ (\text{multiplicity }2sk).
\]
The annihilation-parent compression has eigenvalues
\[
 -\alpha_-s/2\ (\text{multiplicity }(2s+2)k),\qquad
 \alpha_-(s+1)/2\ (\text{multiplicity }2sk).
\]
The ambient eigenvalues are exactly
\[
 \frac{L(L+1)-J(J+1)-s(s+1)}2,
 \qquad L=|J-s|,\ldots,J+s,
\]
with multiplicity $2L+1$. Thus every interlacing bound from these
parents and a single irreducible auxiliary spin is computable with
rational arithmetic.
The accompanying finite scan checks all four endpoint interlacing counts
for $8\le N\le200$ and $1/2\le s\le N+2$. None improves the displayed
combined bound. This finite check is not an all-$N$ or all-auxiliary-space
optimality theorem.
\end{document}
